
\section{Discussion}
\label{sec:discussion}

\textbf{Novelty and construct validity.} ALERT's reusable contribution is the coupling of a point-in-time validity-interval gate (Eq.~\eqref{eq:tpg_score}) with goal-conditioned, coverage-controlled evidence verification---not citators, agents, CSV schemas, or the scoring matrix. The operator is general in form: it applies to any time-stamped, supersession-bearing authority corpus, and we instantiate and evaluate it for law. Table~\ref{tab:temporal_ablation} isolates the mechanism, and the compute-matched baseline shows the gain is not extra LLM calls alone. Because severity labels are triage proxies and 847 labels are retained pre-labels, the human-adjudicated subset is the primary validity result. The two construct-validity checks are separated: external-stack transfer (Protocol~P1) remains the main open generality test, while authoritative-citator agreement for the treatment graph (Protocol~P2) is now \emph{measured} (Appendix~\ref{app:p2}): negative-treatment micro-F1 0.81, a 20.4\% missed-closure rate, and a 3.8\% PLRE decision-flip rate after citator correction, converting treatment-graph correctness from argued to measured and leaving cross-stack transfer as future work.

\textbf{Error propagation and system complexity.} The edge extractor is not trusted as an oracle: a negative-treatment edge must pass citation-span anchoring, date consistency, confidence checks, and Goal-Checker coverage before it can close an interval, or it becomes an unresolved atom forcing targeted retrieval or human escalation. The multi-agent loop is thus a reliability mechanism for localizing graph uncertainty, not an engineering embellishment.
